% truncation_and_diagnostics.tex
%
% Purpose:
%   Discuss Rayleigh truncation, light-line intuition, Bloch multiplier
%   convergence, and practical diagnostics for the global matching matrix.
%
% Main algorithm:
%   No single algorithm is introduced.  The section lists checks that should be
%   run when varying M, N_tot, proxy parameters, and mode-selection thresholds.
%
% Based on:
%   notes/theory/scat_formulation.md, notes/theory/mode_selection.md, and
%   notes/theory/scat_formulation2.md.
%
% Main changes:
%   Heuristic truncation rules are explicitly labeled as practical guidance
%   rather than rigorous error estimates.
%
% Numerical goal:
%   Provide a validation checklist for line-defect scattering and eigenvalue
%   scans.

\section{Truncation Parameters and Diagnostics}

\subsection{Rayleigh truncation}

The truncation parameter \(M\) sets the number \(K=2M+1\) of Rayleigh trace
coordinates.  It controls both the resolution of wall Cauchy traces and the
amount of evanescent coupling retained between separated objects and ports.
For large \(|m|\), real \(k\), and fixed \(\beta\),
\[
  \gamma_m\approx \ii|\beta_m|,
  \qquad
  |\beta_m|\sim\frac{2\pi|m|}{d}.
\]
An evanescent component crossing an open distance \(\ell\) has the heuristic
decay factor
\[
  \ee^{-\operatorname{Im}\gamma_m\ell}.
\]
If each inclusion is separated from its nearest cell wall by at least
\(\delta\), then the gap between adjacent inclusions across a port is roughly at
least \(2\delta\).  A practical high-order coupling heuristic is therefore
\[
  \ee^{-2\operatorname{Im}\gamma_m\delta}.
\]
This is useful intuition for choosing \(M\), but it is not a rigorous error
estimate for the full periodic lead problem.

\subsection{Light-line and periodic dispersion}

For a homogeneous Rayleigh exterior,
\[
  \gamma_m^2=k^2-\beta_m^2 .
\]
Thus the Rayleigh order \(m\) decays in an open direction exactly when
\[
  k<|\beta_m|
\]
for real positive \(k\).  In the simple one-dimensional prototype with
\(0\le\beta\le\pi/d\),
\[
  \text{all Rayleigh orders decay}
  \iff
  k<\min_m|\beta_m|=\beta .
\]

A periodic lead is different.  Its admissibility is selected by the Bloch
multiplier in
\[
  \Asc(k,\beta)z=\lambda\Bsc(k,\beta)z,
\]
not by a scalar Rayleigh light-cone formula.  The Rayleigh light-line remains a
useful diagnostic for the homogeneous background channels, but the periodic
lead dispersion and the selected Cauchy subspace are governed by the truncated
Bloch problem.

\subsection{Diagnostics}

The Rayleigh tail alone does not determine the error in the Bloch multipliers
\(\lambda_j\).  Changing \(M\) changes the truncated scattering matrix and
therefore the generalized eigenproblem itself.  Convergence should be checked
for multipliers, outgoing subspaces, global matrix singular values, and fields.

Useful diagnostics include:
\begin{itemize}
  \item The selected counts \(K_{\mathrm{out}}^\pm\) should be near \(K\) in
  non-degenerate band-gap or regularized cases.
  \item The combined modal matrix
  \[
    \Phi_{\mathrm{all}}^\pm
    =
    \begin{bmatrix}
    D_{\mathrm{in}}^\pm & D_{\mathrm{out}}^\pm\\
    N_{\mathrm{in}}^\pm & N_{\mathrm{out}}^\pm
    \end{bmatrix}
  \]
  should have reasonable numerical rank and conditioning.
  \item Principal angles or projector differences between \(E_{\mathrm{out}}\)
  at successive \(M\) values should decrease in stable parameter regimes.
  \item \(\sigmin(\Atr(k,\beta))\) should converge under refinement of
  \(N_{\mathrm{tot}}\), \(M\), proxy parameters, and selection thresholds.
  \item Residuals of the BIE row and all port matching rows should be checked
  after solving.
  \item Field plots and point evaluations should converge under the same
  refinements, especially near the ports and material interfaces.
\end{itemize}

Potential warning signs include sudden changes in \(K_{\mathrm{out}}^\pm\),
large generalized eigenpair residuals, nearly defective Bloch pencils, grazing
Rayleigh orders \(\gamma_m\approx0\), and strong sensitivity of the field to the
mode-selection tolerance.
